% QUESTAO
Sejam duas matrizes $A$ e $B$, de mesma ordem. Ao adicionarmos os elementos correspondentes (de mesma posição) das matrizes $A$ e $B$, estamos adicionando essas matrizes, ou seja, calculando $A+B$. Assim, dada a matriz $A = \left[ \begin{array}{ccc}
3 & 2 & -1 \\
5 & 0 & 4 \\
2 & -3 & 1
\end{array} \right]$. Determine a soma A+A.

% ALTERNATIVAS
$A+A = \left[ \begin{array}{ccc} 6 & 4 & -2 \\ 10 & 0 & 8 \\ 4 & -6 & 2 \end{array} \right]$
$A+A = \left[ \begin{array}{ccc} 6 & 4 & -2 \\ 10 & 0 & 8 \end{array} \right]$
$A+A = \left[ \begin{array}{ccc} 3 & 2 & -1 \\ 5 & 0 & 4 \\ 2 & -3 & 1 \end{array} \right]$
$A+A = \left[ \begin{array}{ccc} 6 & 4 & 2 \\ 10 & 0 & 8 \\ 4 & 6 & 2 \end{array} \right]$
$A+A = \left[ \begin{array}{ccc} 6 & 4 & -2 \\ 4 & -6 & 2 \end{array} \right]$


